%! Describing the Distribution of a Quantitative Variable — Unit 1, Lesson 1.4 — Homework
% GRADUAL RELEASE: this is the lesson's SCORED practice and the LAST component in the packet.
% Paper homework is the DEFAULT for this course; DeltaMath is an occasional override that
% replaces this page and strikes its score cell on the cover. Started in the last eight minutes
% of the period and finished at home. The AP Practice component that precedes it stays
% assigned-but-unscored.
% THIRD CONTEXT: the notes teach on checkout lines and three histograms, the AP Practice page
% runs on a commute record, and these items are on fresh data again — recall is not the skill.
% \answerspace{H}{} reserves exactly H in the blank; the key passes the answer as the second
% argument, so the two files paginate identically by construction. Keep the macro and every
% \boxguard byte-identical in the blank and the key.
%
% Exercises 2 and 3 reuse displays B and D from exercise 1 on purpose: the context is already
% loaded, and it keeps this component to two pages.
% Exercise 2 IS the target misconception, head-on.
% THE DISPLAYS
%   A symmetric, unimodal:      1 3 6 8 6 3 1
%   B skewed right (pile left): 9 7 4 2 1 1 1
%   C skewed left  (pile right):1 1 2 4 7 9
%   D bimodal with a gap:       2 6 3 0 0 3 6 2
%   Exercise 4 dotplot, sugar (g) in 24 cereals: 2,3,3,4,4,4,5,5,5,5,6,6,6,7,7,8,8,9,10,11,
%   12,13,20,21  -> median 6 g, right tail, two cereals stranded at 20 and 21 g.
\documentclass[12pt]{article}
\usepackage{apstats-article}
\usepackage{apstats-boxes}
\newcommand{\answerspace}[2]{\par\nopagebreak\noindent\begin{minipage}[t][#1][t]{\linewidth}%
  \color{keyred}\bfseries #2\end{minipage}\par}

\begin{document}

\pageheader{Unit 1, Lesson 1.4}{Homework}
% NAMESTRIP: no \namedateperiod here — the name row lives on the cover only.

\vspace{0.10in}

% Authored BYTE-IDENTICALLY in homework/ and homework_key/ — this box is what keeps the
% blank and the key the same number of pages.
\boxguard[8]
\begin{remindbox}
\small
\textbf{This is your graded homework.} It is scored and due next class. You will start it in the
last minutes of the period, so ask about anything that stalls you before you leave, then finish
it on your own.
\end{remindbox}

\vspace{0.12in}

% ── The scored practice — five exercises spanning Topic 1.6 ────────────────
\boxguard[14]
\begin{notesbox}{Describing a Quantitative Distribution}
Write every description \emph{in the context of the problem}, naming the variable and its
units --- that is what earns credit on the AP exam, and it is the standard here.
\begin{enumerate}[leftmargin=1.9em, itemsep=5pt]
  \item \textbf{Match each display to its description.} Write the letter of the display on the
        line beside the description that fits it.

        \vspace{4pt}
        \begin{center}
        \begin{tikzpicture}[scale=0.75]
          % A — symmetric, unimodal
          \foreach \i/\c in {0/1, 1/3, 2/6, 3/8, 4/6, 5/3, 6/1} {
            \fill[navy] ({\i*0.4},0) rectangle ({(\i+1)*0.4},{\c*0.15});
            \draw[white, very thin] ({\i*0.4},0) rectangle ({(\i+1)*0.4},{\c*0.15});}
          \draw[->, thick] (0,0) -- (0,1.8);
          \draw[->, thick] (0,0) -- (3.1,0);
          \node[font=\bfseries\small] at (1.4,-0.55) {A};
          % B — skewed right (pile on the left)
          \begin{scope}[xshift=3.8cm]
            \foreach \i/\c in {0/9, 1/7, 2/4, 3/2, 4/1, 5/1, 6/1} {
              \fill[goldacc] ({\i*0.4},0) rectangle ({(\i+1)*0.4},{\c*0.15});
              \draw[white, very thin] ({\i*0.4},0) rectangle ({(\i+1)*0.4},{\c*0.15});}
            \draw[->, thick] (0,0) -- (0,1.8);
            \draw[->, thick] (0,0) -- (3.1,0);
            \node[font=\bfseries\small] at (1.4,-0.55) {B};
          \end{scope}
          % C — skewed left (pile on the right)
          \begin{scope}[xshift=7.6cm]
            \foreach \i/\c in {0/1, 1/1, 2/2, 3/4, 4/7, 5/9} {
              \fill[navy] ({\i*0.45},0) rectangle ({(\i+1)*0.45},{\c*0.15});
              \draw[white, very thin] ({\i*0.45},0) rectangle ({(\i+1)*0.45},{\c*0.15});}
            \draw[->, thick] (0,0) -- (0,1.8);
            \draw[->, thick] (0,0) -- (3.1,0);
            \node[font=\bfseries\small] at (1.4,-0.55) {C};
          \end{scope}
          % D — bimodal with a gap
          \begin{scope}[xshift=11.4cm]
            \foreach \i/\c in {0/2, 1/6, 2/3, 3/0, 4/0, 5/3, 6/6, 7/2} {
              \fill[goldacc] ({\i*0.35},0) rectangle ({(\i+1)*0.35},{\c*0.15});
              \draw[white, very thin] ({\i*0.35},0) rectangle ({(\i+1)*0.35},{\c*0.15});}
            \draw[->, thick] (0,0) -- (0,1.8);
            \draw[->, thick] (0,0) -- (3.1,0);
            \node[font=\bfseries\small] at (1.4,-0.55) {D};
          \end{scope}
        \end{tikzpicture}
        \end{center}
        \vspace{4pt}

        \begin{enumerate}[label=\textbf{(\roman*)}, leftmargin=2.4em, itemsep=3pt]
          \item Roughly symmetric, with a single peak in the middle. \blank{1.6cm}
          \item Skewed to the right: most values are small, and a few are much larger.
                \blank{1.6cm}
          \item Skewed to the left: most values are large, and a few are much smaller.
                \blank{1.6cm}
          \item Bimodal: two clear peaks with a gap between them. \blank{1.6cm}
        \end{enumerate}

  \item A student looked at \textbf{Display B} and wrote: \emph{``This distribution is skewed to
        the left, because almost all of the data is piled up on the left side.''} The student has
        the direction backwards. Correct the statement, and explain what actually decides the
        name.
        \answerspace{2.4cm}{}
  \item \textbf{Display D} shows the ages, in years, of the 22 people at one community swim
        session. Name the unusual feature this display shows, and give one plausible reason the
        group splits the way it does.
        \answerspace{2.2cm}{}
  \item The dotplot below shows the grams of sugar per serving in each of 24 breakfast cereals.

        \vspace{3pt}
        \begin{center}
        \begin{tikzpicture}[scale=0.88]
          \draw[->, thick] (-0.2,0) -- (5.9,0);
          \foreach \x/\lab in {0/0, 1/4, 2/8, 3/12, 4/16, 5/20} {
            \draw (\x,-0.07) -- (\x,0.07); \node[below, font=\tiny] at (\x,-0.1) {\lab};}
          \foreach \x/\n in {0.5/1, 0.75/2, 1.0/3, 1.25/4, 1.5/3, 1.75/2, 2.0/2, 2.25/1,
                             2.5/1, 2.75/1, 3.0/1, 3.25/1, 5.0/1, 5.25/1} {
            \foreach \k in {1,...,\n} { \fill[navy] (\x,{0.18*\k}) circle (0.06); }}
          \node[font=\tiny] at (2.8,-0.85) {Sugar per serving (grams)};
        \end{tikzpicture}
        \end{center}
        \vspace{2pt}

        Write a complete description of this distribution: all four parts, in context, with the
        variable and its units named.
        \answerspace{2.8cm}{}
  \item \textbf{AP-style multiple choice.} A histogram of the salaries of all 250 employees at
        one company is strongly skewed to the right, with three salaries far above the rest.
        Which statement is best supported? Circle one, then justify it in a sentence.
        \begin{enumerate}[label=\textbf{(\Alph*)}, leftmargin=2.2em, itemsep=2pt]
          \item The distribution is skewed to the left, because most of the salaries are on the
                left side of the graph.
          \item Most of these 250 employees earn near the lower end, while a few earn much more.
          \item Salaries at every company in this industry are skewed to the right.
          \item The three highest salaries should be deleted before the distribution is
                described.
          \item The histogram must be wrong, because a distribution cannot have a long tail.
        \end{enumerate}
        \answerspace{2.0cm}{}
\end{enumerate}
\end{notesbox}

\vspace{0.12in}

% The next-lesson preview closes the packet, so it lives here rather than in AP Practice.
\boxguard[10]
\begin{spiralbox}
\small
\textbf{Next lesson: Summary Statistics for a Quantitative Variable.} Today you said ``about
4 minutes'' and ``most of them between 3 and 9.'' Next time we stop saying \emph{about} and put
actual numbers on the middle of a distribution and on how spread out it is --- and we find out
which of those numbers one strange value can drag around, and which one barely notices it.
\end{spiralbox}

\end{document}
